% !PRINTED none
% Section-opening page: no page number is visible in the running head of any of
% the three bands. Page carries the Abschnitt title, the \S.\,63 sub-heading, the
% opening paragraphs of the section, and a footnote citing the literature.

\gesperrt{Neunter Abschnitt.}

\gesperrt{Congruenzgruppen.}

\S.\,63.

\gesperrt{Functionen-Congruenzen.}

Aus den linearen Substitutionen, deren Bildung und Zusammensetzung wir in den früheren Abschnitten kennen gelernt haben, lassen sich noch eine ganz andere Art endlicher Gruppen ableiten, die \gesperrt{Congruenzgruppen}.
% [sic] Weber prints the plural verb "lassen" with the singular head "eine ... Art"; the source wording is preserved.

Diese Congruenzgruppen haben ein mannigfaches Interesse. Sie stammen zunächst aus der Theorie der elliptischen Functionen, wo sie sich als \gesperrt{Galois}'sche Gruppen der Modulargleichungen einstellen. Sie sind aber auch für die allgemeine Gruppentheorie von Wichtigkeit, weil sie uns ein Mittel geben, ganze Reihen von \gesperrt{einfachen Gruppen} zu bilden. Dies ist um so bemerkenswerther, als, wie wir im vierten Abschnitte gesehen haben, die einfachen Gruppen, wenigstens unter den niedrigeren Gradzahlen, sehr selten sind.

Die Theorie dieser Congruenzgruppen wird nicht nur ausserordentlich verallgemeinert, sondern die herrschenden Gesetze treten weit schärfer hervor, wenn man sich auf eine von \gesperrt{Gauss} herrührende Erweiterung des Congruenzbegriffes stützt, die seitdem mehrfach für die Probleme der Gruppentheorie, besonders von \gesperrt{Galois}, ausgebildet und angewandt worden ist\footnote{\gesperrt{Galois}, \glqq Sur la théorie des nombres\grqq{}. Bulletin des sciences math. de Ferussac. 1830. (Vergl. die Note zu \S.\,151 des ersten Bandes.) --- \gesperrt{Schönemann}, Grundzüge einer allgemeinen Theorie der höheren Congruenzen, deren Modulus eine reelle Primzahl ist. (Crelle's Journ. f. Mathematik, Bd. 31, 1846.) --- \gesperrt{Dedekind}, Abriss einer Theorie der höheren Congruenzen in Bezug auf einen reellen Primzahlmodulus. (Crelle's Journ. f. Mathematik, Bd. 54, 1856.)}.
